\documentclass[12pt,a4paper]{article}
\usepackage[utf8]{inputenc}
\usepackage[indonesian]{babel}
\usepackage{amsmath}
\usepackage{amsfonts}
\usepackage{amssymb}
\usepackage{geometry}

\geometry{margin=2.5cm}

\title{Penyelesaian Titik Berat Kurva Homogen}
\author{Ahmad Fakhrul Bawani}
\date{}

\begin{document}

\maketitle

\section{Soal}
Tentukan titik berat kurva homogen yang dibatasi oleh persamaan:
\begin{equation}
  y^2 = 20x \quad \text{dan} \quad x^2 = 20y
\end{equation}

\section{Penyelesaian}

\subsection{Titik Potong Kurva}
Untuk mencari batas integrasi, kita samakan kedua persamaan:
\begin{align*}
  x^2 &= 20y \implies y = \frac{x^2}{20} \\
  \left( \frac{x^2}{20} \right)^2 &= 20x \\
  \frac{x^4}{400} &= 20x \\
  x^4 - 8000x &= 0 \\
  x(x^3 - 8000) &= 0
\end{align*}
Diperoleh akar-akar:
\begin{itemize}
  \item $x_1 = 0 \implies y_1 = 0$
  \item $x_2 = \sqrt[3]{8000} = 20 \implies y_2 = 20$
\end{itemize}
Titik potong berada di $(0,0)$ dan $(20,20)$.

\subsection{Luas Daerah ($A$)}
Luas daerah di antara kurva atas $f(x) = \sqrt{20x}$ dan kurva bawah $g(x) = \frac{x^2}{20}$ adalah:
\begin{align*}
  A &= \int_{0}^{20} \left( \sqrt{20}x^{1/2} - \frac{x^2}{20} \right) dx \\
  A &= \left[ \frac{2}{3} \cdot \sqrt{20} \cdot x^{3/2} - \frac{x^3}{60} \right]_0^{20} \\
  A &= \left( \frac{2}{3} \sqrt{20} (20\sqrt{20}) - \frac{8000}{60} \right) - 0 \\
  A &= \frac{800}{3} - \frac{400}{3} = \frac{400}{3}
\end{align*}

\subsection{Koordinat Titik Berat $(\bar{x}, \bar{y})$}
Karena kurva simetris terhadap garis $y = x$, maka $\bar{x} = \bar{y}$. Mari kita hitung $\bar{x}$:
\begin{align*}
  \bar{x} &= \frac{1}{A} \int_{0}^{20} x \left( \sqrt{20}x^{1/2} - \frac{x^2}{20} \right) dx \\
  \bar{x} &= \frac{3}{400} \int_{0}^{20} \left( \sqrt{20}x^{3/2} - \frac{x^3}{20} \right) dx \\
  \bar{x} &= \frac{3}{400} \left[ \frac{2}{5} \sqrt{20} x^{5/2} - \frac{x^4}{80} \right]_0^{20} \\
  \bar{x} &= \frac{3}{400} \left( \frac{2}{5} \cdot 8000 - \frac{160000}{80} \right) \\
  \bar{x} &= \frac{3}{400} (3200 - 2000) = \frac{3}{400} (1200) = 9
\end{align*}

Berdasarkan sifat simetri:
\begin{equation*}
  \bar{y} = \bar{x} = 9
\end{equation*}

\section{Kesimpulan}
Titik berat kurva homogen tersebut adalah $(\bar{x}, \bar{y}) = \mathbf{(9, 9)}$.

\end{document}